% =====================================================================
% Decoder-depth ablation of the load-reconstruction surrogate
% ---------------------------------------------------------------------
% Drop-in for Sec. 4.1 (subsec:training_results), subsubsection
% "Network depth": replaces the placeholder paragraph, KEEPING the
% existing \subsubsection{Network depth} heading in chapter3.tex.
% Numbers are the regularized production trainer (Tikhonov
% alpha_reg = 3e-4, tail-compressing output, best-validation selection,
% six-signal input), two seeds {0,100}; identical to the finalized loads
% model except for the decoder depth. Source table:
% recon/analysis/loads_depth_reg_table.csv
% (columns depth_L, seed, d_s, NRMSE_Fy, NRMSE_Mz, NRMSE_combined).
% =====================================================================

The final architectural degree of freedom is the decoder depth $L$, the
total number of hidden layers, split between $\mathcal{NN}_\mathrm{dyn}$ and
$\mathcal{NN}_\mathrm{rec}$ in the ratio $1{:}2$; the baseline retained
elsewhere in this work is $L=6$, with two dynamics and four reconstruction
layers. With the six-signal input fixed, $L$ is swept over $\{6,12,24\}$ at
the two latent dimensions that bracket the sweep of \cref{tab:latent_sweep},
$d_s\in\{1,10\}$, each configuration trained twice from independent seeds with
the regularized trainer described above and the combined NRMSE of
\cref{eq:acc_metrics}. \Cref{tab:depth_loads} reports the results. Doubling the
baseline depth to $L=12$ lowers the combined NRMSE by about $20\%$ at $d_s=1$
and, over the better of the two $L=6$ seeds, by about $30\%$ at $d_s=10$, the
most accurate model being $L=12$ at $d_s=10$ with a combined NRMSE of
$9\cdot10^{-4}$. The added depth also removes the initialisation sensitivity of
the higher-capacity model: at $L=6$ the $d_s=10$ error swings from
$1.35\cdot10^{-3}$ to $6.17\cdot10^{-3}$ between the two seeds, whereas at
$L=12$ both seeds settle near $9\cdot10^{-4}$; at $d_s=1$ the reconstruction is
already seed-robust at the baseline depth. Beyond $L=12$ the training
collapses: at $L=24$ the combined NRMSE rises to $5$--$22\cdot10^{-3}$ for every
configuration and remains seed-dependent, the deep $\tanh$ decoder not being
reliably trainable by the L-BFGS phase at this depth. The production depth
$L=6$ is therefore a conservative choice, and a moderate increase to $L=12$
would improve the load reconstruction; the baseline depth is retained for the
closed-loop model of \cref{subsec:mpc_results}.

\begin{table}[htbp]
\centering
\caption{Open-loop test-set reconstruction accuracy of the aerodynamic loads
($F_y$, $M_z$; combined NRMSE, $\times10^{-3}$) as a function of the decoder
depth $L$ and the latent dimension $d_s$, at the fixed six-signal input, for two
training seeds $\{0,100\}$. $L=6$ is the baseline depth used elsewhere in this
work; its entries are the two-seed values of that architecture and need not
coincide with the best-of-runs selection of \cref{tab:latent_sweep}.}
\label{tab:depth_loads}
\begin{tabular}{cc ccc}
\hline
\rowcolor{bluePoli!40}
$d_s$ & seed & $L=6$ & $L=12$ & $L=24$ \T\B \\
\hline\hline
$1$  & $0$   & $3.66$ & $2.92$ & $9.53$ \T\B \\
$1$  & $100$ & $3.68$ & $3.07$ & $5.18$ \T\B \\
\hline
$10$ & $0$   & $1.35$ & $0.949$ & $22.2$ \T\B \\
$10$ & $100$ & $6.17$ & $0.895$ & $9.18$ \B \\
\hline
\end{tabular}
\end{table}
